\documentclass[11pt]{article}
\usepackage{coursenotes}
\begin{document}
\handoutheader{H3}{Sequential Decisions}{Dynamic treatment regimes, Q-learning, G-estimation and learning from bandit logs}{Meeting 6 (Allocation and policy learning)}

\begin{decision}
A platform sends a coupon in week 1 and decides in week 2 whether to send another, depending on whether the
customer ordered. The question is which \emph{rule} for the two weeks earns the most, not simply whether a coupon works.
The week-1 response is both a result of the first coupon and a reason for the second, which breaks ordinary regression
adjustment.
\end{decision}

\begin{goals}
  \item define a dynamic treatment regime and its value;
  \item state sequential exchangeability and identify a regime's value with the sequential g-formula;
  \item explain why adjusting for an intermediate response biases the effect of the first decision;
  \item find the optimal regime by Q-learning (backward induction) and describe G-estimation of blip functions;
  \item evaluate a policy from bandit logs with inverse-probability and doubly robust estimators.
\end{goals}

%=====================================================================================
\sectionone

\subsection{Regimes and their value}

At decision times $t = 1, 2$ a unit has history $H_1 = X_1$ and $H_2 = (X_1, A_1, X_2)$, where $A_t$ is the action and
$X_2$ is information revealed after the first action (did the customer order?). The outcome $Y$ comes at the end. A
\emph{dynamic treatment regime} is a pair of rules $\mathbf d = (d_1, d_2)$ with $A_1 = d_1(H_1)$ and $A_2 = d_2(H_2)$: the
second action may depend on the response to the first. Its value is $V(\mathbf d) = \E[Y(\mathbf d)]$, where $Y(\mathbf d)$ is the
outcome if the unit were treated according to $\mathbf d$. Static regimes (``coupon both weeks'') are special cases.

\subsection{Identification}

\begin{assum}{Sequential exchangeability and positivity}{seqex}
For every regime, $\{Y(\mathbf d)\} \indep A_1 \mid H_1$ and $\{Y(\mathbf d)\} \indep A_2 \mid H_2$; and every action a regime prescribes
has positive probability given the history at which it is prescribed.
\end{assum}

A \emph{sequential multiple-assignment randomised trial} (SMART), which re-randomises at each decision point, satisfies the
assumption by design. In logs, it requires that every input to each decision be recorded.

\begin{prop}{The sequential g-formula}{gformula}
Under consistency and Assumption~\ref{asm:seqex}, with $a_1 = d_1(x_1)$ and $a_2 = d_2(x_1, a_1, x_2)$,
\[
  V(\mathbf d) = \sum_{x_1}\sum_{x_2} \E[Y \mid x_1, a_1, x_2, a_2]\;\Pr(x_2 \mid x_1, a_1)\;\Pr(x_1).
\]
\end{prop}
\begin{proof}
Work backwards. By exchangeability at $t = 2$ and consistency,
$\E[Y(\mathbf d) \mid x_1, a_1, x_2] = \E[Y \mid x_1, a_1, x_2, a_2]$. Average over $x_2$ given $(x_1, a_1)$: by exchangeability at
$t = 1$, the distribution of $X_2$ among units that received $a_1$ is the distribution it would have under $a_1$. Average
over $x_1$.
\end{proof}

\begin{prop}{Adjusting for the intermediate response is wrong}{mediator}
In a regression of $Y$ on $A_1$, $X_2$ and $A_2$, the coefficient on $A_1$ is not the effect of $A_1$ on $Y$ whenever $A_1$
affects $X_2$ and $X_2$ affects $Y$.
\end{prop}
\begin{proof}
$X_2$ lies on a causal path $A_1 \to X_2 \to Y$; conditioning on it removes that part of the effect (Meeting~1) and, if
anything unrecorded affects both $X_2$ and $Y$, opens a collider path. Omitting $X_2$ instead leaves $A_2$ confounded,
since $A_2$ depends on $X_2$. Neither regression is right; the g-formula averages over $X_2$ instead of conditioning on it.
\end{proof}

\subsection{Finding the best regime: Q-learning}

Define the stage-2 \emph{Q-function} $Q_2(h_2, a_2) = \E[Y \mid H_2 = h_2, A_2 = a_2]$, the best stage-2 rule
$d_2^{\text{opt}}(h_2) = \arg\max_{a_2} Q_2(h_2, a_2)$, and its value $V_2(h_2) = \max_{a_2} Q_2(h_2, a_2)$. Then
$Q_1(h_1, a_1) = \E[V_2(H_2) \mid H_1 = h_1, A_1 = a_1]$ and $d_1^{\text{opt}}(h_1) = \arg\max_{a_1} Q_1(h_1, a_1)$.

\begin{prop}{Q-learning finds the optimal regime}{qlearn}
Under consistency and Assumption~\ref{asm:seqex}, the regime $(d_1^{\text{opt}}, d_2^{\text{opt}})$ maximises $V(\mathbf d)$ over all
regimes.
\end{prop}
\begin{proof}
Dynamic programming. By Result~\ref{prop:gformula}, $V(\mathbf d)$ averages $Q_2$ evaluated at $d_2$'s choice; choosing the maximum at
every $h_2$ can only raise it, whatever $d_1$ is. Given that stage-2 choice, $V(\mathbf d)$ averages $Q_1$ at $d_1$'s choice;
choosing the maximum at every $h_1$ again can only raise it.
\end{proof}

Q-learning fits regression models for $Q_2$ and $Q_1$ backwards. Its weakness is that each Q-function models the whole
outcome, including baseline variation unrelated to the decision, so misspecification can distort the comparison between
actions.

\subsection{G-estimation and structural nested mean models}

A \emph{structural nested mean model} models only the effect of each action, holding later actions at a reference value. The
stage-$t$ \emph{blip} is $\gamma_t(h_t, a_t) = \E[Y(\bar a_{t-1}, a_t, \underline 0) - Y(\bar a_{t-1}, 0, \underline 0) \mid H_t = h_t, A_t = a_t]$, the
effect of $a_t$ followed by the reference action thereafter. \emph{G-estimation} chooses the blip parameters $\psi$ so that the
blipped-down outcome $Y - \sum_{s \ge t}\gamma_s(H_s, A_s; \psi)$, which estimates what $Y$ would have been with the reference
action from $t$ on, is uncorrelated with $A_t - \Pr(A_t = 1 \mid H_t)$ given $H_t$. Under sequential exchangeability this
estimating equation is unbiased if \emph{either} the propensity model \emph{or} a model for the blipped-down outcome's mean is
correct: the dynamic analogue of double robustness. The optimal rule treats when the blip is positive (net of cost), and
models only what matters for the decision.

\subsection{Learning from bandit logs}

A contextual bandit chooses an action $A_i$ for each unit with a logged probability $p_i = \Pr(A_i \mid X_i)$ and records the
reward $R_i$. To evaluate a different policy $\pi$ from these logs:

\begin{prop}{Off-policy evaluation}{ope}
If every action $\pi$ would take has $p_i > 0$, then $\hat V_{\text{IPW}}(\pi) = \frac1n\sum_i \ind{A_i = \pi(X_i)}R_i/p_i$ is unbiased for
$V(\pi)$. The doubly robust estimator
$\hat V_{\text{DR}}(\pi) = \frac1n\sum_i\big[\hat\mu(X_i, \pi(X_i)) + \ind{A_i = \pi(X_i)}(R_i - \hat\mu(X_i, A_i))/p_i\big]$ is unbiased if the logged
probabilities are correct, whatever the reward model $\hat\mu$, and usually far less variable.
\end{prop}
\begin{proof}
Given $X_i$, $\E[\ind{A_i = \pi(X_i)}R_i/p_i \mid X_i] = \E[R(\pi(X_i)) \mid X_i]$ because the action was randomised with probability
$p_i$ (Meeting~5). For the DR estimator, the correction term has conditional mean $\E[R(\pi(X_i)) \mid X_i] - \hat\mu(X_i, \pi(X_i))$.
\end{proof}

Two operational rules follow. Never let an exploring policy drive any action's probability to zero where you may later want
to evaluate it, and log the probabilities, not just the actions. When data are sequential (each customer's later context
depends on earlier actions), these estimators extend with products of probabilities over time, whose variance grows quickly
with the horizon.

\begin{pitfall}[sequential decisions]
Evaluating a two-stage rule with a regression that controls for the week-1 response; learning from logs that recorded
actions but not their probabilities; letting a bandit stop exploring an arm that a later policy will need.
\end{pitfall}

%=====================================================================================
\sectiontwo

\paragraph{Setting.} QuickBite runs a SMART on lapsed customers. In week 1 a coupon is sent at random to half ($A_1$). Customers
who order in week 1 ($X_2 = 1$) get nothing further; among those who do not, a second coupon is sent at random to half ($A_2$).
The outcome $Y$ is orders in weeks 2--4. Each order earns ¥10; each coupon sent costs ¥5 in expectation. Results:
\begin{center}
\begin{tabular}{lcccc}
\toprule
& Ordered in & Responders: & \multicolumn{2}{c}{Non-responders: $\E[Y]$} \\
\cmidrule(lr){4-5}
& week 1 & $\E[Y]$ & 2nd coupon & no 2nd coupon \\
\midrule
Week-1 coupon ($A_1 = 1$) & 0.40 & 1.5 & 0.9 & 0.5 \\
No week-1 coupon ($A_1 = 0$) & 0.25 & 1.6 & 0.9 & 0.3 \\
\bottomrule
\end{tabular}
\end{center}

\paragraph{Step 1: stage 2.} For non-responders, a second coupon adds $0.4$ orders after a week-1 coupon (worth
$10 \times 0.4 - 5 = -$¥1) and $0.6$ orders after none ($+$¥1). So $d_2^{\text{opt}}$: send the week-2 coupon only to non-responders who
did not get one in week 1. The best stage-2 action depends on the history.

\paragraph{Step 2: stage 1.} With week-2 behaviour set optimally, value per customer (week-1 order counts as one order):
\begin{align*}
  Q_1(A_1 = 1) &= 0.40 \times 10(1 + 1.5) + 0.60 \times 10 \times 0.5 - 5 = 10 + 3 - 5 = 8.0, \\
  Q_1(A_1 = 0) &= 0.25 \times 10(1 + 1.6) + 0.75 \times (10 \times 0.9 - 5) = 6.5 + 3 = 9.5 .
\end{align*}
The optimal regime: no coupon in week 1; a coupon in week 2 if the customer has not ordered. Value ¥9.50 per customer.

\paragraph{Step 3: compared with static rules.} By Result~\ref{prop:gformula}, ``coupon both weeks'' earns ¥7.40, ``week 1 only'' ¥8.00,
``never'' ¥8.75, and the optimal regime ¥9.50. ``Coupon both weeks'' generates the most orders ($1.54$ per customer) and the least
profit. Against ``never'', the week-1 coupon adds $1.30 - 0.875 = 0.425$ orders per customer, worth ¥4.25, for a coupon that costs
¥5; and a second coupon after a first one adds less than it costs (Step 1). Responders who got the week-1 coupon do order less
afterwards ($1.5$ versus $1.6$), but that comparison conditions on responding, which the coupon itself changes
(Result~\ref{prop:mediator}).

\begin{exercise}[computation]
Verify the value of ``coupon both weeks'' (¥7.40) with the g-formula, and the number of orders per customer ($1.54$).
\end{exercise}
\answer{Value: $0.40 \times 10 \times 2.5 + 0.60 \times (10 \times 0.9 - 5) - 5 = 10 + 2.4 - 5 = 7.4$. Orders:
$0.40 \times 2.5 + 0.60 \times 0.9 = 1.0 + 0.54 = 1.54$.}

\begin{exercise}[judgement]
An analyst estimates the week-1 coupon's effect by regressing weeks 2--4 orders on $A_1$, $A_2$ and the week-1 order indicator,
Explain what the coefficient on $A_1$ measures and why it should not be read as the coupon's effect.
\end{exercise}
\answer{Conditioning on the week-1 order (a result of $A_1$) compares customers with the same response, removing the coupon's
effect through its main channel (inducing a week-1 order) and leaving a weighted mix of within-group gaps ($-0.1$ among
responders, $+0.2$ and $0.0$ among non-responders without and with a second coupon), which can take either sign. By Result~\ref{prop:mediator} neither including nor excluding $X_2$ gives the effect; use the g-formula
(Result~\ref{prop:gformula}) to compare regimes.}

%=====================================================================================
\sectionthree

\subsection*{Annotated reading}
\begin{readinglist}
\reading{Murphy (2003), ``Optimal dynamic treatment regimes'', \emph{JRSS B}}{The framework, sequential exchangeability and
backward induction}{Sections 1--4}{3}
\reading{Robins (2004), ``Optimal structural nested models for optimal sequential decisions'', \emph{Proceedings of the Second
Seattle Symposium in Biostatistics}}{SNMMs and G-estimation for optimal regimes}{Sections 1--3}{3}
\reading{Chakraborty and Moodie (2013), \emph{Statistical Methods for Dynamic Treatment Regimes}, Springer}{Q-learning,
G-estimation and SMARTs at textbook pace}{Ch.\ 2--5}{2}
\reading{Hern\'an and Robins, \emph{Causal Inference: What If}}{The g-formula and time-varying confounding in the course's
notation}{Ch.\ 19--21}{2}
\reading{Dud\'ik, Langford and Li (2011), ``Doubly robust policy evaluation and learning'', \emph{ICML}}{Off-policy evaluation
from bandit logs}{Sections 1--4}{2}
\reading{Murphy (2005), ``An experimental design for the development of adaptive treatment strategies'', \emph{Statistics in
Medicine}}{How to design a SMART}{All}{2}
\end{readinglist}

\end{document}
